\section{Reproducibility}
\label{sec:repro}

\paragraph{Repository contents.} The repository (link and commit hash on page 1) contains the three notebooks
(\texttt{notebooks/}), the parameter files (\texttt{configs/q*.yaml}), the environment files (\texttt{environment.yml},
\texttt{requirements.txt}, plus \texttt{requirements.lock.txt} with the exact package versions of the reported run), the scripts,
and \texttt{results/}. \texttt{results/} holds the saved split, every CSV/JSON table, every figure, the executed notebooks
(the final logs) and \texttt{environment\_used.json} for each question. Trained checkpoints ($\approx$1.2\,GB) are not committed;
the notebooks regenerate them.

\paragraph{Exact commands.} On a Linux machine with one NVIDIA GPU:
\begin{small}
\begin{verbatim}
git clone <repo> && cd <repo>
git checkout <commit>
conda env create -f environment.yml
conda activate cifar10-study
bash scripts/run_all.sh <ROLL_NUMBER>
python scripts/check_results.py
python scripts/build_report.py --name "<name>" \
   --roll "<roll>" --repo "<url>"
cd report && latexmk -pdf main.tex
\end{verbatim}
\end{small}
\texttt{run\_all.sh} executes Q1, Q2 and Q3 in that order with \texttt{papermill}, injecting \texttt{configs/q*.yaml} and the roll
number. Q2 and Q3 reuse Q1's cached runs when they are present. The reported run was executed on Google Colab
(one notebook per session, the settings in \texttt{configs/}); Q2 loaded its 20 Q1 runs from the cache, while Q3
retrained its classifier (Sec.~\ref{sec:q3}). Q1 was run with the first revision of its notebook, which differs from
\texttt{notebooks/} only in lacking the final cell that records the environment; that cell was run by hand at the end
of the session (see \texttt{results/executed/}). On Colab, each notebook can be run directly, with output folders on Google Drive (see
\texttt{README.md}). \texttt{scripts/make\_split.py} regenerates the split on its own and checks it against the saved file.

\begin{table}[t]
\centering\footnotesize\setlength{\tabcolsep}{2.5pt}
\caption{Every setting needed to reproduce the reported numbers (also in \texttt{configs/}).}
\label{tab:hparams}
\begin{tabular}{@{}l>{\raggedright\arraybackslash}p{0.66\linewidth}@{}}
\toprule
Split & 45k/5k stratified, \texttt{random\_state}=\res{splitstate} \\
Classifiers & AdamW, LR $\in$\{3e-4,1e-3,3e-3\} (val), wd 0.05, cosine, 30 ep., batch 128, fp16, 3 seeds \\
Q1(b) & MLP-XL, 40\% symmetric noise, AdamW 1e-3 const., wd 0, 150 ep.; eval 4$\times$/epoch for 10 epochs, then every epoch \\
Q2(c) & tile sizes 16/8/4, per-image orders seeded by (roll, $t$, split); LR re-selected per condition \\
AEs & $d\in$\{16,128,1024\}, AdamW 1e-3, wd 1e-4, cosine, 30 ep., batch 128, 2 seeds \\
DAE & Gaussian $\sigma$=0.2, noise redrawn per batch; blur $\sigma\in$\{0.5,\dots,2\} (val MSE) \\
Probes & linear: AdamW 2e-3, 40 ep., batch 512, wd$\in$\{0,0.01,0.1,0.5\}; $k$-NN: $k\in$\{10,50\} \\
\bottomrule
\end{tabular}
\end{table}

\begin{table}[t]
\centering\footnotesize\setlength{\tabcolsep}{2.5pt}
\caption{Where every reported number comes from (under \texttt{results/}).}
\label{tab:files}
\begin{tabular}{@{}l>{\raggedright\arraybackslash}p{0.7\linewidth}@{}}
\toprule
Tab.~\ref{tab:models} & \texttt{q1\_outputs/table\_a1\_model\_specs.csv} \\
Tabs.~\ref{tab:q1res}--\ref{tab:q1ev} & \texttt{q1\_outputs/table\_a3\_*.csv}, \texttt{table\_final\_test\_q1a.csv}, \texttt{results\_summary.json} \\
Tabs.~\ref{tab:q1b}--\ref{tab:q1btest} & \texttt{q1\_outputs/table\_b1\_*.csv}, \texttt{table\_b2\_*.csv}, \texttt{table\_final\_test\_q1b.csv} \\
Tabs.~\ref{tab:q2models}--\ref{tab:q2conf} & \texttt{q2\_outputs/table\_q2\_*.csv}, \texttt{results\_summary.json} \\
Tabs.~\ref{tab:q3a}--\ref{tab:q3c} & \texttt{q3\_outputs/table\_q3*.csv}, \texttt{results\_summary.json} \\
Figures & \texttt{q*\_outputs/fig\_*.png} (made by the notebook cells named in the captions) \\
Text in Results & \texttt{q1\_outputs/interpretation.txt}, \texttt{q2/q3\_outputs/conclusion.txt} \\
\bottomrule
\end{tabular}
\end{table}

\paragraph{Determinism.} Seeds fix the split, the initialisations, the batch orders, the permutations and all noise. Remaining
run-to-run variation comes from non-deterministic cuDNN kernels and fp16 accumulation, so a rerun on different hardware
reproduces the numbers to within the reported seed spread, not bit-for-bit. The reported values are the ones in
\texttt{results/}.

\paragraph{Limitations.} (i) No augmentation and short budgets, so absolute accuracies are below standard CIFAR-10 numbers.
This is deliberate, to keep the comparisons controlled. (ii) Three seeds for classifiers and two for autoencoders, which
give wide confidence intervals. (iii) Conclusions about the CNN refer to this global-average-pooled architecture.
(iv) Tile shuffling creates artificial edges at tile borders. (v) The fixed classifier in Q3(b) was trained on sharp
images, so its accuracy on reconstructions mixes information loss with distribution shift; the ceiling rows separate the two.
(vi) Linear and $k$-NN probes measure \emph{accessible} class information only.
